% Standalone CV: all formatting is included for the built-in LaTeX editor.
% The original resume.cls is retained beside this file for reference.
\documentclass[11pt,letterpaper]{article}
\usepackage[T1]{fontenc}
\usepackage[utf8]{inputenc}
\usepackage[margin=0.65in]{geometry}
\usepackage{lmodern}
\usepackage{enumitem}
\usepackage{titlesec}
\usepackage{fancyhdr}
\usepackage{xcolor}
\usepackage{hyperref}
\definecolor{linkblue}{HTML}{184C73}
\hypersetup{colorlinks=true,urlcolor=linkblue,linkcolor=linkblue,
  pdftitle={Duo Zhang -- Curriculum Vitae},pdfauthor={Duo Zhang}}
\setlength{\parindent}{0pt}
\setlength{\parskip}{3pt}
\setlength{\emergencystretch}{2em}
\titleformat{\section}{\large\bfseries}{}{0pt}{}[\titlerule]
\titlespacing*{\section}{0pt}{12pt}{6pt}
\setlist[itemize]{leftmargin=1.25em,topsep=2pt,itemsep=2pt,parsep=0pt}
\setlist[enumerate]{leftmargin=1.7em,label={[\arabic*]},ref=\arabic*,topsep=3pt,itemsep=5pt,parsep=0pt}
\pagestyle{fancy}
\fancyhf{}
\renewcommand{\headrulewidth}{0pt}
\fancyfoot[L]{\small Duo Zhang \quad | \quad Curriculum Vitae}
\fancyfoot[R]{\small \thepage}
\setlength{\footskip}{24pt}
\newcommand{\entry}[2]{\textbf{#1}\hfill #2\par}
\newcommand{\project}[2]{\item \textbf{#1}\enspace #2}
\newcommand{\paperref}[1]{\hyperref[pub:#1]{[\ref*{pub:#1}]}}

\begin{document}
\begin{center}
  {\LARGE\bfseries DUO ZHANG}\\[5pt]
  Ph.D. Candidate in Computer Science, Rutgers University\\[3pt]
  \href{mailto:duo.zhang@rutgers.edu}{duo.zhang@rutgers.edu}
  \enspace $\vert$ \enspace (+1) 917-609-7561
  \enspace $\vert$ \enspace
  \href{https://duozhangrobotics.github.io/}{Personal Website}\\[2pt]
  \href{https://github.com/DuoZhangRobotics}{GitHub}
  \enspace $\vert$ \enspace
  \href{https://scholar.google.com/citations?user=bmEZlbkAAAAJ}{Google Scholar}
\end{center}

\section*{Education}
\entry{Rutgers University}{September 2023--Present}
Ph.D. Candidate, Computer Science. Advisor: Prof. Jingjin Yu.

\entry{New York University}{January 2021--December 2022}
M.S., Courant Institute of Mathematical Sciences.

\entry{Shandong University}{September 2016--June 2020}
B.E., Computer Science (June 2018--June 2020); prior undergraduate study in
Energy and Power Engineering (September 2016--June 2018).
\\Selected honor: National Scholarship, Ministry of Education of China.

\section*{Publications and Preprints}
\textbf{Conference and journal publications}
\begin{enumerate}
  \item\label{pub:sdar} \textbf{Duo Zhang}, Junshan Huang, and Jingjin Yu.
  ``High-Performance Dual-Arm Task and Motion Planning for Tabletop
  Rearrangement.'' \textit{IEEE International Conference on Robotics and
  Automation (ICRA)}, 2026.
  \href{https://arxiv.org/abs/2512.08206}{Paper}.

  \item\label{pub:rvg} \textbf{Duo Zhang}, Zihe Ye, and Jingjin Yu.
  ``Asymptotically-Optimal Multi-Query Path Planning for a Polygonal Robot.''
  \textit{IEEE International Conference on Robotics and Automation (ICRA)}, 2025.
  \href{https://arxiv.org/abs/2409.03920}{Paper}.

  \item\label{pub:sip} \textbf{Duo Zhang}, Chen Liang, Xifeng Gao, Kui Wu, and Zherong Pan.
  ``Provably Feasible Semi-Infinite Program Under Collision Constraints via
  Subdivision.'' \textit{IEEE Transactions on Robotics (T-RO)}, 2024.
  \href{https://ieeexplore.ieee.org/abstract/document/10505800}{Paper}.

  \item\label{pub:grasp} Zherong Pan, \textbf{Duo Zhang}, Changhe Tu, and Xifeng Gao.
  ``Planning of Power Grasps Using Infinite Program Under Complementary
  Constraints.'' \textit{IEEE Robotics and Automation Letters (RA-L)},
  7(1):650--657, 2022.
  \href{https://doi.org/10.1109/LRA.2021.3130376}{Paper}.
\end{enumerate}

\textbf{Preprints}
\begin{enumerate}[resume]
  \item\label{pub:drvg} \textbf{Duo Zhang}, Hechen Zhang, Junshan Huang, and Jingjin Yu.
  ``dRVG: Quadtree-Guided, Resolution-Complete Online Motion Planning for
  Polygonal Robots in Unknown Environments.''
  \textit{arXiv preprint}, 2026.
  \href{https://arxiv.org/abs/2609.31412}{arXiv:2609.31412}.

  \item\label{pub:stprrtc} \textbf{Duo Zhang}, Jintong Li, Junshan Huang, and Jingjin Yu.
  ``ST-pRRTC: Parallel Space-Time RRT-C with Adaptive Goal-Time Forests.''
  \textit{arXiv preprint}, 2026.
  \href{https://arxiv.org/abs/2609.30533}{arXiv:2609.30533}.

  \item\label{pub:favor} \textbf{Duo Zhang}, Zhizhuo Zhang, Xiaoli Bai, and Jingjin Yu.
  ``Efficient B\'ezier Velocity Optimization for Free-Floating Space Manipulators.''
  \textit{arXiv preprint}, 2026.
  \href{https://arxiv.org/abs/2609.31880}{arXiv:2609.31880}.
\end{enumerate}

\newpage
\section*{Research Experience}
\entry{Rutgers University, Algorithmic Robotics \& Control Lab}{September 2023--Present}
Graduate Assistant \enspace $\vert$ \enspace Advisor: Prof. Jingjin Yu

\begin{itemize}
  \project{Geometric and online motion planning.}{Developed rotation-stacked
  visibility graphs for multi-query polygonal-robot planning (RVG), then extended
  the approach to unknown environments through local-roadmap merging and
  quadtree-guided sensing (dRVG). Evaluated against sampling-based planners and
  demonstrated online navigation on a microMVP robot.
  \paperref{rvg}, \paperref{drvg}}

  \project{GPU-accelerated manipulation and space-time planning.}{Developed
  dependency-driven dual-arm rearrangement planning (SDAR) and adaptive goal-time
  forests for planning among moving obstacles (ST-pRRTC). Validated the methods
  on dual UR5e arms and a UR5e operating among Crazyflie quadrotors, respectively.
  \paperref{sdar}, \paperref{stprrtc}}

  \project{Free-floating space manipulation.}{Developed B\'ezier joint-velocity
  optimization with analytical recursive sensitivities (FAVOR) for reaching,
  tracking, and prescribed-time interception under coupled spacecraft--arm motion.
  Evaluated in simulation with a seven-DoF arm and five spacecraft models.
  \paperref{favor}}
\end{itemize}

\entry{Tencent America, Lightspeed \& Quantum Studio}{June 2022--June 2023}
Research \& Development Intern, Copernicus Research Group \hfill Bellevue, WA
\begin{itemize}
  \project{Collision-constrained trajectory optimization.}{Developed a
  semi-infinite programming solver using conservative motion bounds and adaptive
  subdivision to maintain collision-free trajectories; evaluated on articulated
  robot arms and UAVs. \paperref{sip}}
\end{itemize}

\entry{New York University, Courant Institute}{September 2021--December 2022}
Research Assistant with Profs. Daniele Panozzo and Lerrel Pinto \hfill New York, NY
\begin{itemize}
  \project{Deformable-object manipulation.}{Built OpenAI Gym
  environments and PolyFEM tutorials for reinforcement
  learning with deformable objects; set up camera tracking and validated robot
  experiments, including a silicone-ball slingshot.}
\end{itemize}

\entry{Shandong University, Interdisciplinary Research Center}{January--December 2020}
Research Assistant, Prof. Changhe Tu's lab \hfill Qingdao, China
\begin{itemize}
  \project{Power grasp planning.}{Worked with Xifeng Gao and Zherong Pan on
  infinite programming with complementary constraints, kernel-integral relaxation,
  and fast multipole evaluation for optimization-based grasp planning.
  \paperref{grasp}}
\end{itemize}

\entry{University of California, Los Angeles}{July--September 2019}
Research Assistant, Prof. Sriram Sankararaman's lab \hfill Los Angeles, CA\\
Cross-disciplinary Scholars in Science and Technology Program
\begin{itemize}
  \project{Cell-type-specific gene expression prediction.}{Developed
  tensor-component-analysis methods for cell-type-specific
  gene-expression prediction and transcriptome-wide association studies;
  implemented and published an R package for the project.}
\end{itemize}

\section*{Teaching Experience}
\textbf{Teaching Assistant, Rutgers University}
\begin{itemize}
  \item Software Methodology \hfill Fall 2024, Spring 2025
  \item Data Management for Data Science \hfill Spring 2024
  \item Systems Programming \hfill Fall 2023
\end{itemize}

\end{document}
